%% annex/p08/modal-completeness-supplement.tex
\section{Canonical completeness for the four basic normal systems}

\begin{annexitem}{annex:p8.modal.1}{thm:p08-modal-canonical-completeness}{Canonical strong completeness for K, KT, S4, and S5}
\begin{proof}
Fix one of the four Hilbert systems \(L\).  Because necessitation in a
derivation from premises may be applied only to a separately proved theorem,
the ordinary deduction theorem holds: induction on a finite premise derivation
shows that \(\Gamma,A\derives_L B\) implies
\(\Gamma\derives_L A\to B\).  Hence, if
\(\Gamma\not\derives_L A\), the set
\(\Gamma\cup\{\neg A\}\) is \(L\)-consistent.  Enumerate the formulas and,
at each stage, add the next formula when consistency is preserved and otherwise
add its negation.  The union is a maximal \(L\)-consistent set extending
\(\Gamma\cup\{\neg A\}\).  This is the Lindenbaum extension used below.

Let \(W_L\) be all maximal \(L\)-consistent sets and define
\[
 wR_Lv
 \quad\Longleftrightarrow\quad
 \{B:\Box B\in w\}\subseteq v,
 \qquad
 V(p)=\{w:p\in w\}.
\]
The modal existence lemma says that if \(\neg\Box B\in w\), there is
\(v\in W_L\) with \(wR_Lv\) and \(\neg B\in v\).  Otherwise the set
\(\{C:\Box C\in w\}\cup\{\neg B\}\) would be inconsistent.  Some finite
conjunction \(C_1\land\cdots\land C_n\) of its boxed antecedents would then
prove \(B\).  Theorem-only necessitation, distribution, and propositional
reasoning would put \(\Box B\) in \(w\), contradicting maximal consistency.
Lindenbaum extension of that consistent set gives the required \(v\).

Induction on formulas proves the truth lemma
\[
 M_L,w\models B\quad\Longleftrightarrow\quad B\in w.
\]
Boolean cases are the standard maximal-consistency cases.  If \(\Box B\in w\),
the definition of \(R_L\) and the induction hypothesis make \(B\) true at
every successor.  If \(\Box B\notin w\), maximality gives
\(\neg\Box B\in w\), and the existence lemma supplies a successor falsifying
\(B\).  This completes the box case.

The canonical frame has the required property in each extension.  In \(KT\),
\(\Box B\to B\) makes \(wR_Lw\).  In \(S4\), if
\(wR_LvR_Lu\) and \(\Box B\in w\), axiom 4 puts
\(\Box\Box B\in w\), hence \(\Box B\in v\) and \(B\in u\); thus the
relation is transitive.  In \(S5\), reflexivity comes from T.  To prove
Euclideanness, suppose \(wR_Lv,wR_Lu\) and \(\Box B\in v\).  If
\(\Box B\notin w\), then \(\Diamond\neg B\in w\); axiom 5 puts
\(\Box\Diamond\neg B\in w\), so \(\Diamond\neg B\in v\), contradicting
\(\Box B\in v\).  Therefore \(\Box B\in w\), and \(wR_Lu\) gives
\(B\in u\).  Hence \(vR_Lu\).  A reflexive Euclidean relation is symmetric
and transitive, so it is an equivalence relation.

Soundness is induction on derivations: the propositional axioms and modus
ponens preserve local validity, K is valid on every frame, theorem-only
necessitation preserves validity, and T, 4, and 5 have the frame properties
just named.  For completeness, if
\(\Gamma\not\derives_L A\), choose the maximal set \(w\) extending
\(\Gamma\cup\{\neg A\}\).  The truth lemma makes every member of \(\Gamma\)
true at \(w\) and \(A\) false there, on a frame in the required class.  This
contrapositive proves strong completeness for local consequence in all four
rows.
\end{proof}
\end{annexitem}
